\chapter{Сон и ночная работа}
\label{ch:sleep}

Пока Рин общается, память работает «на переднем плане»: извлекает, отвечает,
записывает. Но значительная часть работы памяти живого существа происходит
\emph{в простое} — во сне. Эта глава собирает воедино то, что делает
телепортирующий рой в фоне, и объясняет, почему без этого режима память
деградирует.

\begin{intuition}[«Сон» — это режим, а не отдельный процесс]
Строго говоря, сна как отдельной подсистемы нет. Есть \emph{один всегда включённый}
фоновый разведчик (глава~\ref{sec:background}), чьи параметры зависят от свободных
ресурсов $B_{\text{bg}}=B_{\text{total}}(1-\rho)$. «Сон» — это просто его режим при
\emph{низкой} foreground-нагрузке $\rho$: освобождается бюджет, растут доля телепорта
$\alpha$ и широта покрытия. То же самое, что днём делает «любопытство», но с большим
размахом. Всё, что описано ниже, — про этот низко-$\rho$ режим.
\end{intuition}

\section{Зачем вообще сон}

Три задачи не решаются в режиме реального времени и требуют фоновой переработки:

\begin{enumerate}
\item \textbf{Починка того, о чём мы не знаем, что оно сломано.} Ошибочный вывод
мог продиффундировать в память и дать потомков (ползучая вера,
глава~\ref{ch:epistemics}). Мы не знаем, где именно спрятана ошибка, — поэтому
нужна \emph{постоянная случайная инспекция}.
\item \textbf{Достижение затенённых углов.} Дневной обход \emph{локален по вниманию}
и не дотягивается до узлов, оторванных от активного фокуса
(глава~\ref{ch:topology}). Нужен механизм, гарантирующий покрытие.
\item \textbf{Неспешная консолидация и неожиданные ассоциации.} Обобщение эпизодов и
рождение далёких связей лучше делать без давления текущего разговора.
\end{enumerate}

\section{Телепортирующий обходчик и гарантия покрытия}

Ключевой механизм сна — обход в режиме \emph{стохастического покрытия} (сон-политика
из главы~\ref{ch:walker}): с вероятностью $\alpha$ блуждающий агент
\emph{телепортируется} в случайный узел, иначе делает обычный шаг. Это случайное
блуждание с рестартом; его стационарное распределение задаётся \eqref{eq:ppr}.

\begin{principle}[Почему телепорт необходим, а любопытство не заменит]
Любопытство ищет \emph{известную} неопределённость и потому не покидает окрестности
текущего фокуса. Телепорт делает цепь \emph{эргодической}: каждый узел посещается в
матожидании, ни один не заперт в стоке или на острове (глава~\ref{ch:topology}). Это
\emph{доказуемая гарантия покрытия}, которой у смещённого обхода нет. Именно поэтому
починку неизвестных ошибок можно поручить только сну.
\end{principle}

\begin{intuition}[Анти-энтропийный скраб]
В распределённых хранилищах есть зрелый паттерн: \emph{anti-entropy scrub}
(проверка целостности в ZFS; сверка деревьев Меркла в Dynamo/Cassandra). Система не
знает, какая запись повреждена, поэтому непрерывно и случайно перепроверяет всё
хранилище в фоне. Сон Рин — это анти-энтропийный скраб по графу памяти: не «где
сломалось?», а «регулярно проверяем всё понемногу».
\end{intuition}

\section{Что именно делает рой во сне}

Во время сна телепортирующий рой запускает — с разными вероятностями — весь арсенал
операторов, но с более свободными (и одновременно более строгими по последствиям)
правилами:

\begin{itemize}
\item \textbf{Инспекция противоречий (скраб).} Наткнувшись на узел, столкнувшийся с
другим, обходчик понижает уверенность и помечает потомков на пересмотр (ретракция
по обоснованиям, глава~\ref{ch:epistemics}).
\item \textbf{Консолидация.} Скопления похожих эпизодов сворачиваются в обобщения
(глава~\ref{ch:fusion}).
\item \textbf{Свободный fusion.} Случайно со-активированные далёкие узлы дают
неожиданные гипотезы (пример с Платоном) — провизорные, как и всегда.
\item \textbf{Любопытство наружу.} Узлы с флагом \texttt{need\_to\_discover} могут
инициировать внешний запрос (в источник/поиск), чтобы добыть недостающие
различающие посылки.
\end{itemize}

\begin{intuition}[«Сон» — не поэтическая метафора, а режим низкого приоритета]
Строго говоря, «сон» — это просто фоновый режим обхода с высокой долей телепортов и
разрешёнными операциями записи, запускаемый в простое. Мы называем его сном, потому
что функционально он играет роль сна: консолидирует дневной опыт, чистит
противоречия и «видит» неожиданные связи. Но никакой мистики — это тот же рой, что
и днём, с другими параметрами.
\end{intuition}

\section{Разрешение противоречий: молча или через LLM}
\label{sec:conflict}

Скраб находит много противоречий, и звать тяжёлую LLM на \emph{каждое} — непозволительно.
Нужен критерий: когда конфликт разрешается \emph{молча}, силами графа, а когда его
стоит отдать модели.

\paragraph{Сперва — темпоральная нормализация.} Прежде чем вообще объявлять
противоречие: $X$ и $Y$ конфликтуют, только если это \emph{одна сущность, один слот,
один интервал валидности и scope}. «В 2023 мы жили в Грузии» и «сейчас мы живём в
Португалии» — \emph{не} противоречат: это разные интервалы, оба истинны. Нормализуем
время/сущность/слот и лишь потом считаем конфликт.

\paragraph{Детекция — дёшева, но эмбеддинг ей не детектор.} Что $X$ и $Y$
\emph{логически} исключают друг друга — это \emph{не} про дальность в эмбеддинг-
пространстве. Далёкие не противоречат («кот» и «налоговая декларация» далеки, но не
исключают друг друга), близкие тоже («живу в Лиссабоне» и «живу в Португалии» близки,
но одно \emph{влечёт} другое). Настоящие детекторы противоречия: (1)~знаковое ребро
\textsc{contradicts} (поставлено при записи); (2)~\textbf{schema-level конфликт} —
взаимоисключающие значения одного слота при одном времени/scope («живёт в $X$» vs
«живёт в $Y$» для той же сущности и периода) — дёшево и точно; (3)~NLI/claim-детектор
на трудных случаях. Эмбеддинг годится лишь как \emph{сигнал-кандидат} (подсветить
пары, которые стоит проверить), но не как сам детектор.

\paragraph{Жёсткий гейт прежде марджина: ядро личности не ломается контекстом.}
Прежде чем считать марджин, — жёсткий лексикографический гейт. Если разрыв по
\emph{глубине консолидации} велик (например, свежий \texttt{session}-узел против
\texttt{life}-ядра, $\ge N$ слоёв), младший узел \textbf{не может молча убить}
старший — \emph{независимо} от свежести и числа свидетелей. Иначе спам из 50 узлов за
одну сессию (высокий recency и суммарный indep) линейной суммой перевесил бы
фундамент — «взлом личности контекстом». В таком случае марджин обнуляется —
но, \textbf{чтобы гейт не заморозил устаревшее}, конфликт не выбрасывается, а
\emph{открывает протокол ревизии}: помечает старший узел \texttt{contested\_core\_claim}
$+$ \texttt{need\_to\_revise}, накапливает независимые свидетельства, и старший
уступает только через ревизию (метакогниция/LLM), а не спамом. Так «мы живём в
Грузии» однажды честно уступит «мы живём в Португалии» — но протоколом, а не наплывом
за сессию. (Independence-дисконт уже сбивает вес спама — 50 узлов $\ne$ 50
\emph{независимых} свидетелей, — а гейт добавляет жёсткий backstop поверх.)

Строгость гейта \emph{зависит от} \texttt{node\_kind} (гл.~\ref{ch:substrate}):
\texttt{boundary} почти неприкосновенна без явного события изменения; фактический
\texttt{claim} может устареть (и обязан уступать протоколу); \texttt{preference}
дрейфует плавно; \texttt{self\_model} пересматривается, но медленно и через
метакогницию. То есть $\Delta$глубины-консолидации интерпретируется \emph{через тип
узла}, а не одинаково для всех.

\paragraph{Разрешение — по дешёвому марджину (когда гейт пройден).} Считаем скалярный
\emph{марджин} $m(X,Y)$ на уже хранимых величинах:
\[
m(X,Y)=w_1\,|\conf(X)-\conf(Y)|
+w_2\,\text{recency}
+w_3\,\Delta\text{independent-support}
+w_4\,\Delta\text{consolidation-depth},
\]
где: разрыв уверенности; свежесть / темпоральная валидность (новее и про текущее
состояние — весомее); разрыв \emph{независимой} поддержки (у кого больше независимых
подтверждений — и вот тут прямо переиспользуются Bloom/MinHash-сигнатуры из
главы~\ref{ch:propagation}); глубина консолидации (life-слой перевешивает свежий
session, если у того нет overwhelming support).

\begin{principle}[Критерий «молча vs LLM»]
Если $m(X,Y)>\tau$ — разрешаем \textbf{молча}: проигравший \emph{депрецируется}
(не удаляется — Wikidata-ранг-стиль, остаётся для аудита и на случай отката),
каскадно ретрактятся его производные по TMS (гл.~\ref{ch:epistemics}). $O(1)$-ish,
LLM не трогаем. Если $m(X,Y)\le\tau$ (близко/неоднозначно — нужна семантика: «а точно
противоречат или это контекст-зависимо?») — метим оба \texttt{need\_to\_revise} и
\textbf{откладываем в батч} на один LLM-проход. Раз это сон — не к спеху; дорогую
модель зовём только на реально трудные близкие случаи, не на каждый конфликт.
\end{principle}

Это тот же паттерн, что у планировщика (гл.~\ref{ch:runtime}): дешёвый рефлекс
закрывает подавляющее большинство, LLM-делиберация — только на настоящую
двусмысленность. И заметьте: депрецирование, а не удаление, — потому что «проигравшее»
убеждение может позже вернуться, если придёт новое свидетельство.

\section{Скраб не сразу и не обязан быть быстрым}

Важное свойство: сон чинит память \emph{не сразу}. Мы принципиально не знаем, что и
где сломалось, поэтому вместо адресной починки идёт постоянная случайная инспекция.
Ошибочное убеждение может прожить в памяти какое-то время, транслируя потомков, и
быть исправленным лишь когда телепорт случайно приведёт обходчика в нужный угол или
когда придёт внешний сигнал. Это \emph{нормально} и по-человечески: у живого ума
исправление ползучей веры тоже растянуто во времени. Задача сна — гарантировать, что
исправление \emph{в конце концов} случится, а не что оно случится немедленно.

\begin{workedbox}[Ночь после стрима]
Днём Рин вывела гипотезу «зритель $Q$ токсичен» из одной резкой реплики (провизорно,
$\conf$ низкая). Ночью телепорт случайно приводит обходчика в кластер про $Q$:
там десять тёплых эпизодов. Скраб видит противоречие, понижает «токсичен» ещё
сильнее и помечает на пересмотр. Параллельно консолидация сворачивает десять
эпизодов в «$Q$ — свой, иногда резок под настроение». Наутро модель отношения с $Q$
здоровее, чем была вечером, — и это сделал сон, а не новый разговор.
\end{workedbox}

Псевдокод скраба и сон-политики обхода — в приложении~\ref{ch:algorithms}.

\section{Sleep pressure: сон как саморегуляция, а не расписание}
\label{sec:sleeppressure}

До сих пор сон у нас включался \emph{снаружи}: планировщик увидел долгий простой
($\rho\approx0$) и запустил сон-политику. Это верно как \emph{механика}, но бедно как
\emph{модель}: Рин не «клонит в сон», её туда отправляют по таймеру. Правильное
дополнение — эндогенное \emph{давление сна} (sleep pressure): непрерывная
метакогнитивная величина, которая растёт от накопленной внутренней незавершённости и
сама смещает планировщик к обслуживанию.

\begin{principle}[Sleep pressure — это отложенный техобслуживание-долг, а не координата настроения]
Первый соблазн — считать, что «энергичность» из настроения тратится на глубокие ретривы
и диффузию, и её падение зовёт в сон. Это делает mood \emph{сервисным счётчиком
техдолга} — а он должен отвечать за тон, инициативность, склонность к игре, не за
бухгалтерию обслуживания. Настроение может \emph{отражать} сонливость (устал $\Rightarrow$
рассеяннее, просит паузу), но не быть её \emph{источником}. Источник — отдельная величина.
\end{principle}

И эту величину не надо собирать взвешенной суммой семи «долгов» с коэффициентами
$w_1\ldots w_7$: такие веса нечем откалибровать без downstream-сигнала (та же дисциплина,
что везде в этой книге). Есть заземление чище и без магических чисел. У нас фоновое
обслуживание получает бюджет $B_{\text{bg}}=B_{\text{total}}(1-\rho)$
(гл.~\ref{ch:walker}). Значит sleep pressure — это просто \emph{накопленный отложенный
объём обслуживания}, до которого фон не дотянулся:
\begin{equation}
S(t)=\int_0^t\Big(\text{maintenance\_needed}(\tau)-\text{maintenance\_done}(\tau)\Big)\,d\tau,
\qquad \text{maintenance\_done}\propto(1-\rho).
\label{eq:sleeppressure}
\end{equation}
В загруженные периоды (высокий $\rho$) обслуживание не успевает \emph{по построению}, и
долг растёт сам, без подкрученных коэффициентов. «Долги» из красивых списков — это и есть
\emph{размеры и возраст очередей необслуженного}, которые у нас всё равно заведутся:
дефицит покрытия холодных регионов (coverage-deficit, гл.~\ref{ch:walker}), бэклог
непереваренных кандидатов thought-inbox (гл.~\ref{ch:runtime}), неразрешённые
противоречия (раздел~\ref{sec:conflict}), устаревшие PageRank\,/\,центральности\,/\,индексы,
низко-уверенный мусор. Pressure $=$ «сколько накопилось того, до чего фон не дошёл», а не
эстетская сумма.

\begin{principle}[Foreground локален, память глобальна]
Глубинная причина, зачем сон вообще нужен, не «Рин устала как человек», а: днём она живёт
в локальной зоне фокуса, а память — глобальна. Без сна вся память выродится в историю
последних фокусов. Сон — механизм возврата глобальности: холодные регионы, старые связи,
противоречия, дальние мосты, пересборка карты. Это то же самое, что телепорт-эргодичность
и anti-entropy-скраб выше, только сказанное как \emph{потребность}, а не как расписание.
\end{principle}

\subsection{Стадии сна чистят разные долги}

Sleep pressure — один скаляр (плюс, для диагностики, вектор его слагаемых), но снимается
он \emph{стадийно}, и стадии — это точки на нашем же континууме $\rho$, а не новый ящик в
архитектуре (это \emph{вход} в существующий планировщик, гл.~\ref{ch:runtime}, не отдельный
слой):
\begin{itemize}\itemsep2pt
\item \textbf{Idle / micro-sleep} — паузы между репликами: flush proposals, decay
активации, лёгкая проверка противоречий у фокуса.
\item \textbf{Light sleep} — низкая foreground-нагрузка: локальная консолидация свежих
дневных гипотез, dedup, скраб вокруг активных зон. \emph{Может случаться днём} —
sleep pressure не привязан ко времени суток.
\item \textbf{Deep sleep} — долгое тихое окно (ночь — просто \emph{самое дешёвое} такое
окно): глобальный телепорт-скраб холодных регионов, reindex, пересчёт
PageRank\,/\,центральностей, аудит support-skeleton, prune кандидатов.
\item \textbf{Dream} — не новый механизм, а curiosity-политика (гл.~\ref{ch:walker}) с
высоким телепортом и \emph{commit-authority} $\approx0$: далёкие мосты, странные
ассоциации, много кандидатов — но всё \texttt{pending}, под жёстким anti-fusion и
дисконтом независимости. Dream не \emph{чистит}, а \emph{связывает}; и именно поэтому
его выход карантинится сильнее всего (далёкие мосты — в основном шум).
\end{itemize}
Разные стадии снимают разные слагаемые $S$: light — свежий локальный долг, deep —
структурный и coverage-долг, dream — «интеграционный» дефицит (много со-активированного,
но несвязанного). Отдельную переменную «integration pressure» заводить не нужно — это
слагаемое того же $S$, снимаемое dream-стадией.

\begin{pitfall}[Sleep pressure под теми же тормозами, что Reflecting]
Есть соблазнительное слагаемое «аффективная руминация» — и оно опасно ровно как
самоисполняющаяся тревога (гл.~\ref{ch:mood}): детектор «я туманею, надо переработать»
сам себя кормит, а \emph{вербализованная} сонливость («мне надо уйти разложить хвосты»)
рискует стать перформативным тиком. Поэтому детектор sleep pressure гейтируется теми же
предохранителями, что эндогенное внимание Reflecting (гл.~\ref{ch:runtime}): рейт-лимит,
event-driven каденция, expiry. И foreground-срочность \emph{ветирует} сон: высокий $S$ при
живом разговоре\,/\,дедлайне не отправляет Рин спать — максимум она скажет «после этого мне
надо будет переварить».
\end{pitfall}

\begin{pitfall}[Мерить качество, а не скорость]
Соблазн взять триггером «ретрив замедлился, fusion медленнее». На \emph{растущем} графе
замедление — это чаще \emph{рост}, а не долг (конфаунд с размером). Чистые сигналы долга —
\emph{качественные} и от размера графа не зависящие: доля отклонённых proposals, частота
self-correction, доля low-conf\,/\,нерелевантных узлов в паке, переполнение thought-inbox.
Их и брать, а латентность — нет.
\end{pitfall}

\begin{openq}[Строить после write-path: сейчас нечему копиться]
Инварианты этого раздела фиксируем сразу, но \emph{саму} переменную $S$ строить пока рано:
она интегрирует \emph{очереди необслуженного} — proposals, coverage-долг, неразрешённые
противоречия, — а они появляются только когда работают write-path и fusion
(гл.~\ref{ch:runtime}), которых ещё нет. Так что sleep pressure — это фича «когда система
уже накапливает долги», а не первого среза. Инварианты: \emph{(1)} сон не cron, ночь —
лишь дешёвое окно для deep; \emph{(2)} $S$ — метакогнитивная величина из отложенного
обслуживания \eqref{eq:sleeppressure}, не координата mood и не сумма магических весов;
\emph{(3)} дневной инсайт разрешён, сон его проверяет\,/\,интегрирует\,/\,убивает;
\emph{(4)} стадии снимают разные долги; \emph{(5)} детектор сонливости — под тормозами
Reflecting, foreground ветирует.
\end{openq}
